%======================================================================
\section{Uniqueness results}\label{sec:uniqueness}
%======================================================================

\providecommand{\uqpmat}[4]{\left(\begin{smallmatrix}#1&#2\\#3&#4\end{smallmatrix}\right)}%
This section proves uniqueness for admissible pairs whose supports are close to those of $\pz$. Section~\ref{ssec:barrier} shows,
with soft analysis only, how far positivity alone goes. From Section~\ref{ssec:support} on, the input is the Fourier interpolation
theorem of Bondarenko, Radchenko and Seip \cite{BRS}, which does not assume RH. We use the following notation.
$\vartheta(\tau)=\sum_{n\in\ZZ}e^{\pi in^2\tau}$ is Jacobi's theta function and $\Gamma_\vartheta$ the theta group, generated by
$S=\uqpmat{0}{-1}{1}{0}$ and $T^2=\uqpmat1201$ and viewed, as in \cite{BRS}, in $\mathrm{PSL}_2(\RR)$, so that matrices are identified up to
sign and $S^2=1$ (as matrices, $S^2=-I$); we write $\tau=\sigma+iv\in\mathbb H$ and $\omega=e^{2\pi i/3}$. $\boldsymbol\mu$ is the
M\"obius function, $\one_\square(m)=1$ if $m$ is a perfect square and $0$ otherwise, and $\Lambda(\sqrt m):=0$ if $m$ is not a square.
Recall $\BRS=\{\log m/(4\pi):m\ge1\}$; note $\xin n=\log(n^2)/(4\pi)$, so $\xiPP\subset\BRS$.

\subsection{Why arithmetic is needed}\label{ssec:barrier}

How far can two admissible pairs differ, if one uses only positivity and the gap $(-\xin2,\xin2)$ in the prime measures? For a real
measure $\nu$ on $[\xin2,\infty)$ write $\nu^{\mathrm s}:=\nu+\nu(-\,\cdot)$ for its even extension.

\begin{lemma}[Homogeneous directions]\label{lem:homdir}
Let $\mathcal D$ be the real vector space of pairs $\Delta=(\Delta\mu,\Delta\nu)$, with $\Delta\mu$ an even real Radon measure
on $\RR$ and $\Delta\nu$ a real Radon measure on $[\xin2,\infty)$, such that $|\Delta\mu|([T-1,T+1])\ll\log(3+|T|)$,
$|\Delta\nu|([X-1,X+1])\ll e^{\pi X}$, and $\Spec_{\Delta\mu,\Delta\nu}(F)=0$ for every $F\in\Tc$. Then:
\begin{enumerate}[label=\textup{(\alph*)}]
\item $\Kad-\Kad\subseteq\mathcal D$.
\item For $\Delta$ with these growth bounds, $\Spec_\Delta=0$ on $\Tc$ if and only if $\widehat{\Delta\mu}=-\frac1{2\pi}(\Delta\nu)^{\mathrm s}$
in $\mathcal D'(\RR)$; in that case $(\Delta\nu)^{\mathrm s}$ is tempered.
\item Consequently $\widehat{\Delta\mu}=0$ on $(-\xin2,\xin2)$: any two admissible pairs satisfy $\int\varphi\dd\mu_1=\int\varphi\dd\mu_2$
for every even real $\varphi$ with $\widehat\varphi\in C_c^\infty((-\xin2,\xin2))$.
\end{enumerate}
\end{lemma}

The proof is in Appendix~\ref{app:proofs-rigidity}. By (b), the elements of $\mathcal D$ are pairs of measures in Fourier duality with a
spectral gap. Those whose measures are countable sums of point masses are Fourier summation formulas in the sense of
\cite{Gon23,GV25,Ved26}; in general they are not, since a smooth even density on $[\xin2,\infty)$ and its Fourier transform give an element
of $\mathcal D$ with two continuous sides.

\begin{proposition}[Two-sided rigidity]\label{prop:rigidity}
Let $\Delta=(\Delta\mu,\Delta\nu)\in\mathcal D$ and let $Z\subset\RR$ be countable.
\begin{enumerate}[label=\textup{(\alph*)}]
\item If $0\notin Z$, $\Delta\mu\ge0$ off $Z$, $\sum_{z\in Z}(\Delta\mu)_-(\{z\})<\infty$, and $\Delta\nu\ge0$ off a
Lebesgue-null set, then $\Delta=0$.
\item If $0\notin Z$, $\Delta\mu\ge0$ off $Z$ and $(\Delta\nu)_-$ is a finite measure, then $\Delta=0$.
\item If $(\Delta\mu)_-$ is carried by a finite set, then $\Delta=0$.
\end{enumerate}
\end{proposition}

The proof (Appendix~\ref{app:proofs-rigidity}) uses only that every even $\varphi\ge0$ with $\widehat\varphi$ supported in
$(-\xin2,\xin2)$ has the same integral against $(\Delta\mu)_+$ and $(\Delta\mu)_-$, by (c). Applied to $\Delta=(\mu-\muz,\nu-\nuz)$
under RH, with $Z=\Zz$ (which does not contain $0$, since $\zeta(\frac12)\ne0$; indeed $(1-2^{1-\sigma})\zeta(\sigma)=\sum_{n\ge1}(-1)^{n-1}n^{-\sigma}>0$
for $0<\sigma<1$ \cite[(25.2.3)]{DLMF}), it shows that no admissible pair other than $\pz$ arises from $\zeta$'s pair by only adding mass
on one side, or by deleting a finite amount of mass on either side. The hypotheses cannot be weakened, by the following examples.

\begin{proposition}[Poisson examples]
Let $\beta\ge\xin2$.
\begin{enumerate}[label=\textup{(\alph*)}]
\item $\Delta\mu:=\beta\dd t-\sum_{n\in\ZZ}\delta_{n/\beta}$, $\Delta\nu:=2\pi\beta\sum_{m\ge1}\delta_{\beta m}$ defines a non-zero
$\Delta\in\mathcal D$ with $\Delta\nu\ge0$ and $\Delta\mu\ge0$ off $Z=\beta^{-1}\ZZ$, which contains $0$; $(\Delta\mu)_-$ has infinite mass.
\item $\Delta\mu:=\beta\dd t-\sum_{n\in\ZZ}\delta_{(n+1/2)/\beta}$, $\Delta\nu:=2\pi\beta\sum_{m\ge1}(-1)^m\delta_{\beta m}$ defines a non-zero
$\Delta\in\mathcal D$ with $\Delta\mu\ge0$ off $Z=\beta^{-1}(\ZZ+\frac12)$, which does not contain $0$, and $\Delta\nu\ge0$ off the null
set $\{\beta m:m\text{ odd}\}$; both $(\Delta\mu)_-$ and $(\Delta\nu)_-$ have infinite mass.
\end{enumerate}
\end{proposition}

\begin{proof}
By Poisson summation, $\widehat{\sum_n\delta_{n/\beta}}=\beta\sum_m\delta_{\beta m}$ and
$\widehat{\sum_n\delta_{(n+1/2)/\beta}}=\beta\sum_m(-1)^m\delta_{\beta m}$, while $\widehat{\beta\dd t}=\beta\delta_0$. Hence
$\widehat{\Delta\mu}=-\frac1{2\pi}(\Delta\nu)^{\mathrm s}$ in both cases, the growth bounds are clear, $\Delta\nu$ lives on
$[\beta,\infty)\subseteq[\xin2,\infty)$, and Lemma~\ref{lem:homdir}(b) gives $\Delta\in\mathcal D$. The sign statements are read off.
\end{proof}

So the sign and gap conditions of Proposition~\ref{prop:rigidity}, taken alone, do not force homogeneous rigidity. The examples are not
shown to be differences of admissible pairs for $\zeta$'s data, so they do not rule out every soft proof of \condU; they show that such a
proof needs more than these conditions. For $\beta=\xin2$ the exceptional set in (b) lies inside the prime-power nodes, so the obstruction
is on the zero side. The rest of this section supplies arithmetic through the basis of \cite{BRS}.

\subsection{Uniqueness on the support of \texorpdfstring{$\zeta$}{zeta}'s pair}\label{ssec:support}

For $s\in\CC$ let $\mathsf F(\tau,s)=\sum_{N\ge0}\alpha_N(s)e^{\pi iN\tau}$ be the unique $2$-periodic analytic function of
moderate growth on $\mathbb H$ satisfying
\begin{equation}\label{eq:BRS-FE}
\mathsf F(\tau,s)+(\tau/i)^{-1/2}\,\mathsf F(-1/\tau,s)=\mathfrak p_s(\tau):=(\tau/i)^{-s}+(\tau/i)^{s-1/2}
\end{equation}
(principal branches); this is $F^-_{1/2}(\tau,s)$ of \cite[Theorem~3.1]{BRS}. Its coefficients grow polynomially in $N$, are entire
in $s$, and satisfy $\alpha_N(\frac12-s)=\alpha_N(s)$ and $\alpha_N(\bar s)=\overline{\alpha_N(s)}$ (both by uniqueness in
\eqref{eq:BRS-FE}); hence $\alpha_N(s)\in\RR$ on the line $\Re s=\frac14$. Let $b_N$ $(N\ge0)$ be the even Fourier
eigenfunctions with eigenvalue $+1$ of Radchenko and Viazovska \cite{RV19} (the functions $b_N^-$ of \cite[\S3.4.1]{BRS}), so that
$b_N(\sqrt n)=\delta_{Nn}$ for $n\ge1$, and put
\[
B_N(x):=\sum_{k\ge1}b_N(kx)\qquad(x>0).
\]

\begin{definition}[The BRS basis]
By \cite[Theorem~1.1]{BRS}, which assumes neither RH nor simplicity of the zeros, there are even entire functions $U_m$
$(m\ge1)$ and $V_{\rho,j}$ ($\rho$ a non-trivial zero with $\Im\rho>0$, $0\le j<m(\rho)$, where $m(\rho)$ is the multiplicity)
such that
\begin{gather*}
U_m^{(j)}(t_\rho)=0\ \ (j<m(\rho)),\qquad \widehat{U_m}\Big(\frac{\log m'}{4\pi}\Big)=\delta_{mm'},\\
V_{\rho,j}^{(j')}(t_{\rho'})=\delta_{(\rho,j),(\rho',j')},\qquad \widehat{V_{\rho,j}}\Big(\frac{\log m}{4\pi}\Big)=0
\end{gather*}
for all $m,m'\ge1$. They decay rapidly in vertical strips: see the paragraph on the rapid decay of the basis functions in the proof of
\cite[Theorem~1.1]{BRS} (\cite[\S4.3]{BRS}). The bounds used there for $\zeta^*$ and for the coefficients $\alpha^\pm_{m,1/2}$ are uniform
for $\Re s$ in compact sets, so the decay is uniform on every closed strip $\lvert\Im z\rvert\le\frac12+\delta$, and the real and imaginary
parts of these functions lie in $\Tc$. Explicitly \cite[\S\S4.2--4.3]{BRS},
\begin{equation}\label{eq:Um}
U_m(z)=\frac{m^{-1/4}}{2\pi}\,h_m\big(\tfrac12+iz\big),\qquad h_m(s)=\frac{\GR(s)\zeta(s)}2\sum_{d^2\mid m}\boldsymbol\mu(d)\,\alpha_{m/d^2}(s/2),
\end{equation}
and $\widehat{U_m}(\xi)=m^{-1/4}e^{\pi\xi}u_m(e^{2\pi\xi})$ with $u_m(x)=\sum_{d^2\mid m}\boldsymbol\mu(d)B_{m/d^2}(x)$.
We call $U_m$ and $V_{\rho,j}$ the \emph{BRS functions}.
\end{definition}

By Lemma~\ref{lem:explicit-formula}, applied to real and imaginary parts,
\begin{equation}\label{eq:W-BRS}
\Arch(U_m)=\frac1\pi\Lambda(\sqrt m)\,m^{-1/4},\qquad \Arch(V_{\rho,0})=2m(\rho),\qquad \Arch(V_{\rho,j})=0\quad(j\ge1).
\end{equation}
Indeed $U_m$ vanishes at every $t_\rho$ and $\widehat{U_m}(\xin n)=\widehat{U_m}(\log(n^2)/4\pi)=\delta_{m,n^2}$. For $V_{\rho,j}$ the
prime sum vanishes; the zeros with negative imaginary part are the $1-\rho'$ with $\Im\rho'>0$, with $t_{1-\rho'}=-t_{\rho'}$, so
by evenness $\sum_{\rho'}V_{\rho,j}(t_{\rho'})=2\sum_{\Im\rho'>0}m(\rho')V_{\rho,j}(t_{\rho'})=2m(\rho)\delta_{j0}$. These are the values
of \cite[\S4.4]{BRS}, with two slips corrected: there the first value is printed as $\pi^{-1}\Lambda(n)n^{-1/2}$ for a square $m=n$,
whereas \eqref{eq:W-BRS} gives $\pi^{-1}\Lambda(k)k^{-1/2}$ at $m=k^2$, and the second as $2$ instead of $2m(\rho)$.

\begin{proposition}[Uniqueness on $\zeta$'s support]\label{thm:support}
Let $(\mu,\nu)\in\Kad$ with $\mu$ carried by $\Zz$ and $\nu$ carried by $\BRS$. Then RH holds and $(\mu,\nu)=\pz$. Positivity of
$(\mu,\nu)$ is not used beyond making the identity meaningful.
\end{proposition}

This is essentially the observation made in \cite[\S4.4]{BRS}, that the Riemann--Weil functional has the values
\eqref{eq:W-BRS} in the interpolation basis, combined with the uniqueness statement \cite[Corollary~1.1]{BRS}; we record it as a
corollary of \cite[Theorem~1.1]{BRS}.

\begin{proof}
Write $\mu=\sum_{\gamma\in\Zz}w_\gamma\delta_\gamma$ and $\nu=\sum_{m\ge1}\nu(\{\frac{\log m}{4\pi}\})\delta_{\log m/4\pi}$. Every pairing
below is a finite sum, by the interpolation properties. Suppose $\rho_0$ is a zero with $\Im\rho_0>0$ and $\Re\rho_0\ne\frac12$.
Every real ordinate $\gamma\in\Zz$ is $\pm t_\rho$ for an on-line zero $\rho\ne\rho_0$ with $\Im\rho>0$, so $V_{\rho_0,0}(\gamma)=0$;
and $\widehat{V_{\rho_0,0}}$ vanishes on $\BRS$. Hence $\Spec_{\mu,\nu}(V_{\rho_0,0})=0$, whereas $\Arch(V_{\rho_0,0})=2m(\rho_0)\ne0$ by
\eqref{eq:W-BRS}. So RH holds. For $\gamma>0$ in $\Zz$ and $\rho=\frac12+i\gamma$, testing on $V_{\rho,0}$ gives
$w_\gamma+w_{-\gamma}=2m(\rho)$, i.e.\ $w_\gamma=m_\gamma$ since $\mu$ is even; so $\mu=\muz$. Testing on $U_m$ gives
$\frac1\pi\nu(\{\frac{\log m}{4\pi}\})=\frac1\pi\Lambda(\sqrt m)m^{-1/4}$, i.e.\ $\nu=\nuz$.
\end{proof}

The proof uses the identity $\Arch=\Spec_{\mu,\nu}$ only for $F\in\{U_m,V_{\rho,j}\}$, and no positivity. In this form it applies to
real measures $\mu$ (even) and $\nu$ carried by $\Zz$ and $\BRS$, and we use it so in the proof of Theorem~\ref{thm:Gfin}.

For admissible pairs the hypothesis on $\nu$ can be dropped, and an atom of $\mu$ at the origin can be allowed.

\begin{theorem}[Zero-side support]\label{thm:zerosupport}
Let $(\mu,\nu)\in\Kad$, and assume that $\mu$ is carried by $\Zz\cup\{0\}$. Then RH holds and $(\mu,\nu)=\pz$.\astranote
\end{theorem}

No hypothesis on $\nu$ is made beyond admissibility: neither discreteness nor support in $\BRS$. In particular, for admissible pairs,
Proposition~\ref{thm:support} holds without its hypothesis on $\nu$. The proof uses $\nu\ge0$, so the signed form of
Proposition~\ref{thm:support} used in the proof of Theorem~\ref{thm:Gfin} is not strengthened.

\emph{Sketch of the proof} (Appendix~\ref{app:proofs-zerosupport}). Let $\nu^\sharp$ be the image of $\nu$ under $\xi\mapsto x=e^{2\pi\xi}$, a
positive measure on $[2,\infty)$ with $\nu^\sharp([R,2R])\ll R^{1/2}$ by Lemma~\ref{lem:apriori}; so the lattice functional
$L(f):=\int\sqrt x\sum_{k\ge1}f(kx)\dd\nu^\sharp(x)$ is continuous on even Schwartz functions. Testing the admissibility identity on the
functions $U_m$, which vanish on $\Zz$, and inverting over square divisors gives $L(b_N)=\frac12(\log N)\one_\square(N)+2\mu(\{0\})\Xi(0)\alpha_N(\frac14)$.
As the interpolation series of Radchenko and Viazovska converges in Schwartz space, every even Schwartz function $f$ with $\hat f=f$ and
$f(0)=0$ satisfies $L(f)=\sum_{n\ge2}(\log n)f(n)-\frac12\mu(\{0\})\zeta(\frac12)\int_0^\infty f(x)x^{-1/2}\dd x$. The function
$g(t)=\sin^2(\pi t)/(\pi^2(t^2-1))$ vanishes at the integers, is positive at all other points of $(1,\infty)$, and $\hat g$ is supported in
$[-1,1]$ and vanishes at $0$ and $\pm1$. Gaussian regularisations of $g+\hat g$ in this identity give
$0\le\int\sqrt x\sum_kg(kx)\dd\nu^\sharp(x)=-\frac12\mu(\{0\})\zeta(\frac12)\int_0^\infty(g+\hat g)(x)x^{-1/2}\dd x\le0$, as both $\zeta(\frac12)$
and the last integral are negative. Hence $\mu(\{0\})=0$ and $\nu^\sharp$ is carried by the integers $n\ge2$. Testing on $b_{m^2}$ and inverting
over divisors gives $\nu^\sharp(\{m\})=\Lambda(m)m^{-1/2}$, that is, $\nu=\nuz$, and Theorem~\ref{thm:logic}(b) gives RH and $\mu=\muz$.


\subsection{Countably many extra zeros}\label{ssec:Gfin}

Fix a countable symmetric set $E\subset\RR\setminus\Zz$, and put $E_+:=E\cap[0,\infty)$ and $s_e:=\frac14+\frac{ie}2$.

\begin{lemma}[Structure of pairs with extra atoms]\label{lem:structure}
Let $\mu$ be an even real measure carried by $\Zz\cup E$ and $\nu$ a real measure carried by $\BRS$ such that
$\Arch(F)=\Spec_{\mu,\nu}(F)$, with absolutely convergent integrals, for $F=U_m$ $(m\ge1)$ and $F=V_{\rho,j}$. Write $a_e:=\mu(\{e\})$, $w_\gamma:=\mu(\{\gamma\})$,
$\ell_m:=m^{1/4}\nu(\{\frac{\log m}{4\pi}\})$, $C_N:=\sum_{k^2\mid N}\ell_{N/k^2}$,
\[
\tilde\varrho_e:=\frac{a_e\,\Xi(e)}{\frac14+e^2}\ (e>0),\qquad \tilde\varrho_0:=2a_0\,\Xi(0),\qquad c_N(E):=\sum_{e\in E_+}\tilde\varrho_e\,\alpha_N(s_e)\in\RR .
\]
Then the series $c_N(E)$ converge absolutely, and for every $N\ge1$
\begin{equation}\label{eq:CN}
C_N=c_N(E)+\tfrac12(\log N)\one_\square(N).
\end{equation}
Moreover: $c_1=c_2=c_3=0$ when $\nu$ is carried by $[\xin2,\infty)$; $w_\gamma=m_\gamma-\frac12\sum_{e\in E}a_eV_{\rho,0}(e)$
for $\gamma>0$ and $\rho=\frac12+i\gamma$; $\frac12\sum_ea_eV_{\rho,0}(e)=m(\rho)$ for every zero $\rho$ off the line with $\Im\rho>0$;
$\sum_ea_eV_{\rho,j}(e)=0$ for $1\le j<m(\rho)$; and $(\mu,\nu)=\pz$ if and only if $a=0$, and, under the summability condition
\eqref{eq:summable} below, if and only if $c(E)\equiv0$. If $(\mu,\nu)\in\Kad$, then $C_N\ge0$ for all $N$; in particular $c_N(E)\ge0$ for
every non-square $N$.
\end{lemma}

The proof is in Appendix~\ref{app:proofs-structure}.

At a non-square $N$ there is no room: $c_N(E)\ge0$ outright. At squares, $\zeta$'s own prime mass $\frac12\log N$ gives a positive
margin, which is why a soft argument cannot work there.

\begin{theorem}[Countably many extra zeros]\label{thm:Gfin}
Let $E\subset\RR\setminus\Zz$ be countable and symmetric. Let $\mu$ be an even real measure carried by $\Zz\cup E$ and $\nu$ a real
measure carried by $\BRS$ such that $\Arch(F)=\int F\dd\mu+\frac1\pi\int\hF\dd\nu$, with absolutely convergent integrals, for every BRS
function $F$ (in particular, if the identity holds on $\Tc$ with absolutely convergent integrals). Assume $\nu(\{\frac{\log m}{4\pi}\})\ge0$
for every $m\equiv2\pmod3$, $\mu(\{0\})\ge0$ if $0\in E$, and, with $a_e:=\mu(\{e\})$ and
$A(T):=\sum_{e\in E,\,0<e\le T}|a_e|\,|\zeta(\frac12+ie)|$,
\begin{equation}\label{eq:summable}
\sum_{e\in E,\,e>0}|a_e|\,|\zeta(\tfrac12+ie)|\,(1+e)^{-1/2}<\infty,\qquad\text{or, more generally,}\qquad A(T)=o(T^{1/2}).
\end{equation}
Then RH holds, $\mu=\muz$ and $\nu=\nuz$. In particular every admissible pair carried by $(\Zz\cup E)\times\BRS$ and satisfying
\eqref{eq:summable} is $\pz$. For finite $E$, \eqref{eq:summable} holds trivially.\footnote{The extension from finite to countable $E$ (with Lemmas~\ref{lem:aggregate}
and~\ref{lem:meanvalue}) is due to Astra (OpenAI), contributed during an independent review of an earlier version of this paper, and
independently verified. The case of finite $E$, with Lemma~\ref{lem:structure} for finite $E$ and Lemmas~\ref{lem:cusp0}
and~\ref{lem:cusp-class2}, is from that earlier version.}
\end{theorem}


\begin{proof}
By Lemma~\ref{lem:structure}, \eqref{eq:CN} holds. If $N\equiv2\pmod 3$ and $k^2\mid N$, then $3\nmid k$, so
$N/k^2\equiv2\pmod3$; hence $c_N(E)=C_N=\sum_{k^2\mid N}\ell_{N/k^2}\ge0$ for all $N\equiv2\pmod3$ (no such $N$ is a square). Since
$a_0\ge0$ and $\Xi(0)=-\frac18\pi^{-1/4}\Gamma(\frac14)\zeta(\frac12)>0$, we have $\tilde\varrho_0\ge0$. Lemma~\ref{lem:meanvalue} in
Appendix~\ref{app:cusp}, whose hypotheses are \eqref{eq:summable} and these signs, now gives $\tilde\varrho_e=0$ for every $e\in E_+$. As
$\Xi(e)\ne0$ for $e\notin\Zz$, $a\equiv0$, so $\mu$ is carried by $\Zz$, and the proof of Proposition~\ref{thm:support} (which uses no
positivity) gives RH, $\mu=\muz$ and $\nu=\nuz$.
\end{proof}

\emph{The argument in more detail.} Positivity enters only through the residue class $2\bmod3$, and the work is to see the extra
zeros in the coefficients $c_N(E)$ of that class. Their generating function
\[
\Theta_E(\tau):=\sum_{N\ge1}c_N(E)e^{\pi iN\tau}=\sum_{e\in E_+}\tilde\varrho_e\big(\mathsf F(\tau,s_e)-\alpha_0(s_e)\big)
\]
is a combination of the modular integrals of \eqref{eq:BRS-FE}, absolutely convergent under \eqref{eq:summable}, and by orthogonality of
additive characters the class $2\bmod3$ is picked out by three cusps:
\[
\Theta_E^{(2)}(v):=\sum_{N\equiv2\,(3)}c_N(E)e^{-\pi Nv}=\tfrac13\Big[\Theta_E(iv)+\omega\,\Theta_E\big(\tfrac23+iv\big)+\bar\omega\,\Theta_E\big(\tfrac43+iv\big)\Big].
\]
The points $0$, $\frac23$ and $\frac43$ are cusps of $\Gamma_\vartheta$ equivalent to $\infty$. At $0$, \eqref{eq:BRS-FE} itself gives the
behaviour of $\mathsf F(iv,s)$ as $v\to0^+$. At $\frac23$, the word $ST^2ST^2S$ maps $\frac23$ to $\infty$, and unrolling \eqref{eq:BRS-FE}
along it writes $\mathsf F(\frac23+iv,s)$ as $-\alpha_0(s)\vartheta(\frac23+iv)$ plus explicit powers of $v$ and terms that are analytic at
$v=0$; at $\frac43$ one uses $\mathsf F(-\bar\tau,s)=\overline{\mathsf F(\tau,s)}$ (Lemmas~\ref{lem:cusp0} and~\ref{lem:cusp-class2}). Two
cancellations follow. Each cusp contributes a theta term of size $v^{-1/2}$, and these combine to a multiple of
$\sum_{n^2\equiv2\,(3)}e^{-\pi n^2v}$, which vanishes identically because no square is $\equiv2\pmod3$; and the constants
$\alpha_0(s_e)$ cancel because $1+\omega+\bar\omega=0$. For a single $s$ the remaining terms are explicit (Lemma~\ref{lem:cusp-class2}).
For countably many $s_e$ we estimate them in aggregate (Lemma~\ref{lem:aggregate}): \eqref{eq:summable}, Stirling's formula and an
argument margin of size $v$ at the two intermediate points of the word show that, with $u:=\log(1/v)$,
\begin{gather*}
v^{1/4}\,\Theta_E^{(2)}(v)=Q(u)+o(1)\qquad(v\to0^+),\qquad Q(u):=B_0+\sum_{e\in E_+,\,e>0}\big(B_ee^{ieu/2}+\overline{B_e}e^{-ieu/2}\big),\\
B_e=\frac{\tilde\varrho_e}3\big(1-3\cdot9^{-s_e}\big),\qquad B_0=\frac{2\tilde\varrho_0}3\big(1-\sqrt3\big),
\end{gather*}
with $\sum|B_e|<\infty$. Since $|3\cdot9^{-s_e}|=\sqrt3\ne1$, $B_e=0$ only if $\tilde\varrho_e=0$. A mean-value argument concludes
(Lemma~\ref{lem:meanvalue}). The almost periodic function $Q$ has mean $B_0$ and mean square $B_0^2+2\sum_{e>0}|B_e|^2$, and
$\liminf_{u\to\infty}Q(u)\ge0$ because $\Theta_E^{(2)}\ge0$. So $B_0\ge0$; as $B_0\le0$, $B_0=0$; and a bounded function with mean $0$ that is
eventually $\ge-\eta$, for every $\eta>0$, has mean square $0$, so every $B_e$ vanishes. In an earlier version of this paper, for finite
$E$, the last step used Landau's theorem on Dirichlet series with non-negative coefficients \cite[Theorem~1.7]{MV07} instead.

\emph{Related results.} The identities of Lemma~\ref{lem:structure} are the explicit-formula form of a remark in \cite[\S4.1]{BRS}:
by Bochner's converse theorem \cite{Boc51}, a Dirichlet series with the functional equation of the kernel and finitely many poles is
determined up to the kernels themselves, which strengthens Knopp's abundance principle \cite{Kno00}. The restriction to a residue
class is a twist by additive characters, the standard tool of the degree-one classification of the Selberg class
\cite{CG93,KP99,Sou05}; there it is not combined with positivity. The closest uniqueness theorem we found is
Booker's converse theorem for positive $L$-data \cite[Definition~1.3, Theorem~1.7]{Boo15}. There the zero multiplicities are real,
integral outside a finite set, and non-negative outside a finite set, and positive data of degree less than $5/3$ are $1$ or shifted
Dirichlet $L$-functions; the prime side is a priori supported on $\{\log n\}$ with growth conditions on the coefficients. Theorem~\ref{thm:Gfin}
is complementary to it, not a generalisation: it allows non-integral weights at infinitely many ordinates and prime mass at the non-square
points of $\BRS$, imposes no growth condition on the prime side, and uses positivity on the prime side, on one residue class, instead of on the zero side; but
it requires the zero measure to be carried by $\Zz$ and countably many further points satisfying \eqref{eq:summable}, and the prime measure
by $\BRS$. In the language of Fourier summation formulas \cite{Gon23}, the signed
pairs that the BRS basis produces on $\Zz\cup E$ are exactly what the theorem excludes under positivity.


\subsection{The magic-function principle}

\begin{corollary}[Magic-function principle]\label{cor:magic}
Let $F\in\Cone$ with $\Arch(F)=0$, and let $E:=Z(F)\setminus\Zz$ be the set of its real zeros outside $\Zz$.
\begin{enumerate}[label=\textup{(\alph*)}]
\item If $E\subseteq\{0\}$, then \condU\ holds.
\item If $\widehat Z(F)\subseteq\BRS$ and
\begin{equation}\label{eq:summableF}
\sum_{n\ge0}(1+n)^{-1/2}\log(3+n)\,\max_{e\in E\cap[n,n+1]}|\zeta(\tfrac12+ie)|<\infty
\end{equation}
(the maximum over an empty set being $0$; this holds whenever $E$ is finite), then \condU\ holds.
\end{enumerate}
In either case, if \condE\ holds as well, then RH holds and $\Kad=\{\pz\}$.
\end{corollary}

Part~(a) is due to Astra (OpenAI), contributed during an independent review of an earlier version of this paper, and independently
verified; it needs no condition on the zeros of $\hF$.

\begin{proof}
Let $(\mu,\nu)\in\Kad$. As $0=\Arch(F)=\int F\dd\mu+\frac1\pi\int\hF\dd\nu$ with both integrals non-negative, $\int F\dd\mu=0$; so
$\mu(\{F>1/k\})\le k\int F\dd\mu=0$ for every $k$, and $\mu$ is carried by $Z(F)\subseteq\Zz\cup E$. (a) Theorem~\ref{thm:zerosupport} applies. (b) By
Lemma~\ref{lem:compslack}, $\nu$ is carried by $\widehat Z(F)\subseteq\BRS$. The set $E$ is symmetric, because $F$ is even, and countable,
because $F$ is analytic and not identically zero. By Lemma~\ref{lem:apriori}, $\mu(E\cap[n,n+1])\le C\log(3+n)$, so
$\sum_{e\in E,\,e>0}\mu(\{e\})|\zeta(\frac12+ie)|(1+e)^{-1/2}$ is at most $C$ times the sum in \eqref{eq:summableF}, and
Theorem~\ref{thm:Gfin} applies.
\end{proof}

As in the proof that the Leech lattice is the unique optimal lattice in its dimension \cite{CK09}, the zeros of an extremal function
determine the extremal configuration. With extra real zeros, the prime-side condition $\widehat Z(F)\subseteq\BRS$ in (b) is needed by the
method (Remark~\ref{rem:extraprimes}). The converse, that \condU\ implies the existence of such an $F$, is not known: given \condE\ it
amounts to a strict complementarity statement for an infinite-dimensional linear programme (Problem~\ref{op:strict-compl}). The functions of
Section~\ref{sec:family} have the form $\Xi^2H$, for which $Z(\Xi^2H)=\Zz\cup\{\text{real zeros of }H\}$; so (a) applies to $\Xi^2H$
whenever $H$ has no real zeros other than possibly $0$.

\subsection{Integer weights and extra prime atoms}

With integer zero weights, as for an $L$-function, \condU\ becomes a uniqueness problem for Beurling generalised number systems with
Riemann's functional equation. For prime measures on $\{\xin n:n\ge2\}$ this is Hamburger's theorem \cite{Ham21}, and for prime
measures that generate a uniformly discrete multiplicative semigroup, a setting related to Lagarias's study of Beurling integers with the
Delone property \cite{Lag99}, it is treated in \cite{N4} (in preparation) with the theorem of Lev and Olevskii on measures with uniformly discrete
support and spectrum \cite{LO15}; Problem~\ref{op:Beurling} asks for the general case.

\begin{remark}[Extra prime atoms]\label{rem:extraprimes}
If $\mu$ is carried by $\Zz\cup\{0\}$, Theorem~\ref{thm:zerosupport} excludes every extra prime mass: atoms at any position, finitely or
infinitely many, and diffuse mass. With extra zeros, Theorem~\ref{thm:Gfin} still needs the prime measure on $\BRS$, and this is where the
method stops. An extra prime-side atom
at $\eta$ enters the analogue of the identities \eqref{eq:CN} through the numbers $B_N(x_\eta)$, $x_\eta=e^{2\pi\eta}$, which are the Fourier
coefficients of a second modular integral built from $\vartheta(x_\eta^2\tau)$. Its sums over a class of quadratic non-residues
modulo an odd prime have a leading term of order $v^{-1/2}$ at the cusps, given by Gauss sums, when $x_\eta^2$ is rational; when
$x_\eta^2$ is irrational that leading term vanishes. So an atom with $x_\eta^2\notin\QQ$ is invisible to the leading-order argument of
Lemma~\ref{lem:meanvalue}; see \cite{N2}.
\end{remark}
